\section{Теоретични основи}

За да бъде разбрана и аргументирана архитектурата на разглежданата система, е необходимо да бъдат въведени няколко основни теоретични понятия. Макар проектът да има подчертано практическа насоченост, неговата реализация стъпва върху утвърдени концепции от областта на уеб инженерството и софтуерната архитектура: клиент--сървърно взаимодействие, многослойна организация, MVC модел, обектно-релационно моделиране, управление на сесия и валидация на данни. Тази глава не цели изчерпателен теоретичен обзор, а по-скоро въвежда именно онези идеи, които са пряко необходими за разбиране на последващото описание на системата.

\subsection{Клиент--сървърна природа на уеб приложението}

В основата на всяка уеб система стои клиент--сървърният модел. Клиентът, най-често уеб браузър, инициира HTTP заявка към сървъра. Сървърът интерпретира заявката, прилага бизнес логика, извлича или записва данни и връща отговор, който може да бъде HTML документ, пренасочване, файл или структуриран формат като JSON. Тази обща схема е в сила и за разглежданата система: браузърът на потребителя изпраща заявки към контролерите на приложението, а те в зависимост от маршрута връщат страници, извършват пренасочване или генерират REST отговор.

От инженерна гледна точка важна е идеята, че протоколът HTTP е по своята природа безсъстояниен. Това означава, че отделните заявки са независими една от друга, освен ако приложението не въведе собствен механизъм за проследяване на потребителския контекст. В настоящия проект тази роля се поема от \texttt{HttpSession}, където се съхранява електронната поща на текущо удостоверения потребител. Така системата постига достатъчно прост модел на сесийно състояние, подходящ за учебно приложение, без да въвежда по-сложни механизми като токени, филтри и централизирана сигурност.

\subsection{Многослойна архитектура}

Многослойната архитектура е един от най-устойчивите организационни модели в enterprise разработката \cite{fowler,sommerville}. Нейната идея е различните отговорности в системата да бъдат разделени на слоеве с относително ясни граници. В типичния случай могат да бъдат разграничени визуален слой, приложен слой за управление на заявки и потоци, бизнес слой за домейн логика и слой за достъп до данни.

Предимството на този модел е, че промени в един слой могат да бъдат извършвани с по-малко странични ефекти върху останалите. Например промяна в HTML шаблон не би следвало да налага промяна в entity класовете, а промяна в начина на съхранение на данни не би следвало да засяга непосредствено логиката на формите. Разбира се, в реални системи винаги има определено взаимодействие между слоевете, но многослойният модел създава дисциплина, която намалява вероятността от силно преплитане на отговорности.

В настоящия проект многослойността е реализирана чрез контролери, услуги, repository интерфейси и шаблони. Този подход е особено подходящ за образователна разработка, защото позволява лесно да се проследи движението на данните от потребителската форма към persistence слоя и обратно.

\subsection{Моделът MVC}

Model--View--Controller е класически архитектурен шаблон за уеб приложения и графични системи \cite{springmvc}. В контекста на Spring MVC той се проявява по следния начин:

\begin{itemize}[leftmargin=*]
    \item \textbf{Model} представя данните, които се подават към визуализацията;
    \item \textbf{View} е шаблонът, който визуализира тези данни пред потребителя;
    \item \textbf{Controller} обработва входящата заявка и избира коя view компонента да бъде върната.
\end{itemize}

Важно е да се подчертае, че в Spring MVC терминът „модел“ не съвпада напълно с домейн модела на приложението. В конкретната реализация домейн обектите са entity класовете, а обектът \texttt{Model} в контролерите представлява контейнер за атрибути, които ще бъдат достъпни от Thymeleaf шаблона. Това разграничение е съществено, защото показва, че MVC не е заместител на домейн модела, а механизъм за организиране на взаимодействието между потребителската заявка и визуализирания резултат.

В разглежданата система MVC моделът е видим почти във всеки маршрут. Контролерът получава заявка за регистрация, логин, каталог, количка или профил; извиква съответна услуга; поставя резултата в модела; и връща име на шаблон. Тази последователност е една от причините проектът да бъде добър пример за ясна учебна архитектура.

\begin{figure}[H]
    \centering
    \begin{tikzpicture}
        \node[appblock] (browser) {Уеб браузър};
        \node[appblock, below=1.4cm of browser] (controller) {Controller};
        \node[appblock, below left=1.5cm and 2.4cm of controller] (service) {Service слой};
        \node[appblock, below right=1.5cm and 2.4cm of controller] (view) {Thymeleaf view};
        \node[appblock, below=1.4cm of service] (repo) {Repository слой};
        \node[appblock, below=1.4cm of repo] (db) {MySQL база данни};

        \draw[appline] (browser) -- node[right]{HTTP заявка} (controller);
        \draw[appline] (controller) -- node[left]{бизнес логика} (service);
        \draw[appline] (service) -- (repo);
        \draw[appline] (repo) -- (db);
        \draw[appline] (controller) -- node[right]{модел} (view);
        \draw[appline] (view) |- node[pos=0.25,right]{HTML} (browser);
    \end{tikzpicture}
    \caption{Обобщен поток на обработка при сървърно рендирана заявка}
    \label{fig:mvc-request-flow}
\end{figure}

Фигура \ref{fig:mvc-request-flow} илюстрира основния път на заявката в системата. Макар той да изглежда линеен, зад всяка стъпка стои отделен набор от отговорности: контролерите управляват маршрутизацията, услугите формулират правилата на домейна, repository компонентите делегират работата с базата данни, а шаблоните осигуряват потребителската визуализация.

\subsection{Dependency Injection и конфигурация на компоненти}

Една от основните идеи в Spring е \textit{dependency injection}, тоест доставяне на необходимите зависимости от външен контейнер, вместо създаването им на ръка във всеки клас. Практически това позволява класовете да бъдат по-слабо свързани, по-лесно тествани и по-ясно конфигурирани \cite{springboot}. Когато дадена услуга зависи от repository интерфейс, model mapper или password encoder, тя ги получава чрез конструктора си, вместо сама да ги инстанцира.

В текущия проект това е реализирано последователно: контролерите приемат необходимите услуги и инфраструктурни компоненти чрез конструкторна инжекция, а конфигурационният клас дефинира bean инстанции за \texttt{ModelMapper} и \texttt{PasswordEncoder}. Това решение е важно не само за чистотата на кода, но и за архитектурната проследимост: когато се разглежда даден клас, веднага се вижда от кои други компоненти зависи и каква роля изпълнява.

\subsection{Домейн модел и обектно-релационно съответствие}

Информационните системи често трябва да представят едни и същи понятия в два различни свята: обектния свят на програмния език и релационния свят на базата данни. Обектно-релационното съответствие се стреми да преодолее тази разлика чрез механизъм, при който класовете и връзките между тях се свързват с таблици, колони и външни ключове \cite{springdatajpa}. В Java екосистемата това обикновено се реализира чрез JPA анотации върху entity класове.

В разглежданата система класове като \texttt{User}, \texttt{Company}, \texttt{Product}, \texttt{ShoppingCartEntity} и \texttt{HistoryEntity} са маркирани като entity обекти и описват основните релации в домейна. Така се постига сравнително естествена връзка между предметната област и persistence слоя. Например един потребител има множество фирми, количка и история на поръчки; продуктът принадлежи към фирма; а записът в количката свързва конкретен продукт с конкретен потребител и желано количество.

Тук е важно да се направи разграничение между \textit{домейн модел} и \textit{модели за пренос на данни}. В проекта не всички данни се обменят чрез entity обектите. За някои операции са въведени binding модели за входни форми, service модели за вътрешен трансфер между слоеве и view модели за изход към интерфейса. Това е добър архитектурен избор, тъй като ограничава директното изтичане на persistence структурата към потребителския интерфейс.

\subsection{Валидация на входни данни}

Въвеждането на данни от потребител е една от най-честите причини за логически грешки, нарушаване на ограничения в базата или несъответствие между очакваните и действителните стойности. Поради това валидацията на входните данни е фундаментална характеристика на всяка надеждна информационна система. В Spring приложенията това често се реализира чрез Java Bean Validation анотации и обработка на \texttt{BindingResult} \cite{hibernatevalidator}.

В конкретния проект ограниченията са дефинирани основно в binding моделите. За полета като име, адрес, количество, цена и други са зададени правила от типа минимална и максимална дължина, непразна стойност или положително число. Този подход има няколко предимства. Първо, позволява формалното описание на допустимите входни стойности да бъде концентрирано на едно място. Второ, дава възможност контролерът да реагира систематично при невалиден вход, като върне потребителя към съответната форма с индикация за грешка. Трето, намалява риска в service слоя да достигат невалидни обекти.

Въпреки това е важно да се отбележи, че валидацията в текущия проект е частична. Например някои binding модели не съдържат богато множество от ограничения, а част от логическите проверки се извършват на по-късен етап в контролерите или услугите. Това е приемливо за учебен проект, но при по-зряла система би било полезно валидационните правила да бъдат още по-последователни и формализирани.

\subsection{Управление на потребителска идентичност и защита на паролите}

Системите за електронна търговия оперират с чувствителни данни и поради това дори базовият модел на автентикация трябва да бъде проектиран разумно. Един от най-важните принципи е паролите никога да не се съхраняват в явен вид, а да бъдат обработвани чрез устойчив механизъм за хеширане \cite{springsecuritycrypto,owasp}. В текущия проект тази роля се изпълнява от \texttt{Pbkdf2PasswordEncoder}, който осигурява по-добра защита от директно съхраняване или използване на бързи, неподходящи за пароли хеш функции.

Моделът на удостоверяване в проекта е сравнително опростен. При успешен вход електронната поща на потребителя се записва в \texttt{HttpSession}. От тази стойност впоследствие се определя дали дадена заявка е позволена и кой е текущият потребител. Този подход е лесен за разбиране и достатъчен за малка учебна система, но има ограничена изразителност спрямо по-пълноценни решения със сигурностен контекст, роли, филтри и декларативни правила за достъп. По-късно в документа това ще бъде разгледано като естествена възможност за бъдещо надграждане.

\subsection{Сървърно рендирани шаблони и локализация}

Thymeleaf е шаблонен механизъм за сървърно рендиране, който интегрира естествено HTML маркировка и изрази за работа с данни от модела \cite{thymeleaf}. Неговото предимство е, че прави HTML шаблоните сравнително четими и при липса на изпълнение, а в същото време позволява условна визуализация, итериране по колекции, двупосочно свързване на форми и използване на локализирани съобщения.

В настоящия проект Thymeleaf се използва като основен механизъм за изграждане на потребителския интерфейс. Формите за регистрация, вход, добавяне на компания, добавяне на продукт, профил, каталог, количка и история са реализирани именно по този начин. Допълнителен теоретичен аспект тук е локализацията. Чрез message bundles и интерцептор за смяна на езика системата демонстрира базов модел за поддръжка на повече от един език, като конкретно предлага английски и български вариант на част от интерфейсните текстове.

\subsection{REST интерфейс и представяне на данни}

Макар приложението да е доминиращо сървърно рендирано, в него е наличен и примерен REST интерфейс, връщащ информация за максималната цена на продукт. Теоретичното значение на тази част е, че показва как една и съща бизнес логика може да бъде използвана през различни форми на представяне. Вместо HTML, при REST сценария резултатът се сериализира до JSON и може да бъде използван от JavaScript код в браузъра или от външен клиент. Подобен стил на взаимодействие е в съзвучие с идеите на уеб архитектурата и ресурсно ориентираното представяне на данни \cite{fielding}.

Дори този относително малък REST пример има методологична стойност. Той показва, че границата между традиционно MVC приложение и API ориентирано приложение не е абсолютно строга. Една добре структурирана система може да комбинира двата подхода и да предоставя различни изгледи върху една и съща домейн логика според нуждите на клиента.

\subsection{Обобщение на теоретичната рамка}

Разгледаните понятия очертават теоретичния фундамент на настоящата разработка. Клиент--сървърният модел определя начина, по който потребителят взаимодейства със системата. Многослойната архитектура организира вътрешните ѝ отговорности. MVC моделът структурира уеб заявките и визуализациите. Обектно-релационното моделиране свързва домейн обектите с базата данни. Валидацията ограничава невалидния вход. Хеширането на пароли и управлението на сесия осигуряват базова сигурност. Thymeleaf и локализацията оформят представящия слой, а REST крайна точка показва алтернативен начин за експониране на данни.

Тази теоретична рамка не е само фон; тя е пряко материализирана в структурата на проекта. Следващата глава преминава от общите понятия към конкретните изисквания и потребителски сценарии, които системата трябва да удовлетворява.
